% Generated by report/make_restated.py from theory/B3_prop59_tightness.tex: do not edit (edit the note and rerun).
\begin{restated}{Remark}{rem:B3-orders}{B3}{0}{reading the two attribution orders}
The attribution is path dependent, as for any decomposition of a difference of a nonlinear function. Order B (undo from the end) attributes to each inequality its full effect on the weighted contraction; order A keeps the max structure, so that undoing (I4) and (I2), which touch $c_D$ alone, has \emph{no} effect as long as the memory coefficient $c_M$ is the larger one (regime $\kappa<9n/16$ for the textbook constants), and undoing (I3), which touches $c_M$ alone, has no effect as long as $c_D$ is the larger one. Undoing (I1) is different: $r_{\mathrm{sq}}$ enters both coefficients (through $2\gamma^2r_{\mathrm{sq}}\LEff\muEff$ in $c_D$ and $2\gamma^2r_{\mathrm{sq}}r_{\mathrm{var}}/\Delta$ in $c_M$), so the step $\rho^A_3\to\rho^A_4$ lowers both of them and has an effect in both regimes, except in the degenerate cases $r_{\mathrm{sq}}=1$, or $r_{\mathrm{var}}=0$ with $c_M$ the larger coefficient. Reporting both orders therefore shows, in addition to the individual contributions, which component ($D$ or $M$) is the bottleneck of the bound. Since $\rho$ is close to $1$, $\nu_k\approx\delta_k$; per effective pass the natural unit is $n\nu_k$.
\end{restated}
